% Fixtures
\newcommand{\cblock}[2]{
\left\|
\begin{matrix}
#1 & #2\\
-#2 & #1
\end{matrix}
\right\|
}

% Content
\section{Postać normalna macierzy rzeczywistych symetrycznych, antysymetrycznych oraz ortogonalnych}

W tym podrozdziale omówimy postaci normalne dla macierzy rzeczywistych symetrycznych, antysymetrycznych oraz ortogonalnych na podstawie Twierdzenia Spektralnego i pokażemy kilka przykładów. Także omówimy dlaczego przypadek zespolony jest bardziej skomplikowany.

\vspace{0.3em}
\medskip
\noindent\textbf{Forma normalna macierzy symetrycznej}
\vspace{0.3em}


Przypadek operatora symetrycznego, któremu odpowiada macierz symetryczna, jest prosty, ponieważ możemy bezpośrednio skorzystać z Twierdzenia Spektralnego dla operatora hermitowskiego \ref{theorem:hermitian_spectral}.  

Macierz symetryczna jest szczególnym przypadkiem macierzy hermitowskiej w przestrzeni euklidesowej, ponieważ $S = S^* = \overline{S^T}$, redukuje się do $S = S^T$ dla macierzy rzeczywistej. Macierz hermitowska ma tylko rzeczywiste wartości własne. Z tego możemy wywnioskować, że zawsze będzie istnieć baza ortonormalna złożona z rzeczywistych wektorów własnych w której macierz $S$ jest diagonalna, więc \textbf{forma normalna macierzy symetrycznej} $S$ ma następujący wygląd:
\begin{equation}
    S = O\{\lambda_i\}O^T, \quad O^T = O^{-1} 
\end{equation}
gdzie $O$ macierz ortogonalna, której kolumny tworzą bazę ortonormalną złożoną z wektorów własnych macierzy $S$, a $\lambda_i \in 
\R$ wartości własne macierzy $S$.

\vspace{0.3em}
\medskip
\noindent\textbf{Forma normalna macierzy ortogonalnej i antysymetrycznej}
\vspace{0.3em}

Przypadek macierzy ortogonalnej i antysymetrycznej, jest nieco trudniejszy, ponieważ dla tych macierzy wartości własne mogą być liczbami zespolonymi, więc nie możemy wprost skorzystać z Twierdzenia Spektralnego w przestrzeni rzeczywistej. 

Dla badania formy normalnej macierzy ortogonalnej i antysymetrycznej rozszerzmy przestrzeń euklidesową $V$ do przestrzeni unitarnej $\V$.

\begin{enumerate}[leftmargin=1cm, itemsep=2pt, label=(\arabic*)]
    \item Wektory z $V$ będziemy nazywać \textit{rzeczywistymi}, a wektory z $\V$ \textit{zespolonymi}.
    \item W przestrzeni unitarnej $\V$ każdy wektor zespolony $z \in \V$ możemy przedstawić jako $z = x + iy$, gdzie $x, y$ wektory rzeczywiste z $V$.
\end{enumerate}

Niech $A$ operator liniowy w $V$, wtedy w $\V$ operator liniowy $A$ rozszerza się do:
$$ Az = A(x + iy) = Ax + iAy$$
$$ A\bar{z} = A(x - iy) = Ax - iAy$$
Taki operator $A$ będziemy nazywać \textit{rzeczywistym} w przestrzeni unitarnej $\V$.

Niech $A$ operator normalny rzeczywisty w $\V$. Warunek operator normalnego $A^*A = AA^*$ w przypadku rzeczywistym redukuje się do $A^TA = AA^T$. Na mocy Twierdzenia Spektralnego operator normalny w przestrzeni unitarnej $\V$ ma bazę ortonormalną $U$, złożoną z wektorów własnych, w której ma postać diagonalną, gdzie na przekątnej leżą wartości własne $A$. 

Równanie charakterystyczne operatora $A$ będzie miało współczynniki rzeczywiste\cite[s. 22--23]{jeffrey2006}, więc wartości własne zespolone zawsze będą w parze: $\lambda, \bar{\lambda} \in \C$. Z $Az = \lambda z $ wynika $\overline{Az} = \overline{\lambda z} \implies  A\bar{z} = \bar{\lambda}\bar{z}$, tj. sprzężonym wartościom własnym odpowiadają sprzężone wektory własne.

Z bazy ortonormalnej $U$ wybierzmy dwa dowolne wektory: $z, \bar{z}$, które będą odpowiadać parze wartości własnych: $\lambda, \bar{\lambda}$. Z własności bazy ortonormalnej $z, \bar{z}$ są ortogonalne, więc możemy stworzyć z nich dwuwymiarową podprzestrzeń przestrzeni $\V$:
$$ W = \mathrm{lin}(z, \bar{z}) $$
W tej przestrzeni możemy stworzyć bazę rzeczywistą:
\[
\begin{cases}
z = x + iy & \text{gdzie } x, y \in V \\
\bar{z} = x - iy & \text{gdzie } x, y \in V
\end{cases}
\]

$$ x = \frac{1}{2}(z + \bar{z}), \quad y = \frac{1}{2i}(z - \bar{z})$$
Wektory $x, y \in V$ są ortogonalne:
$$ \inner{z}{\bar{z}} = 0 \implies \inner{x + iy}{x - iy} = 0 \implies \inner{x}{x} + 2i\inner{x}{y} - \inner{y}{y} = 0 $$
Ostatnią równość możemy przepisać jako:
$$\norm{x}^2 + 2i\inner{x}{y} - \norm{y}^2 = 0$$

Norma wektora jest zawsze liczbą rzeczywistą oraz $\inner{x}{y}$ też jest liczbą rzeczywistą ponieważ $x, y$ należą do przestrzeni euklidesowej $V$. Więc:
\begin{equation}\label{eq:real_vec_orthogonal}
    \inner{x}{y} = 0 \implies x \perp y
\end{equation}
Zauważmy, że:
$$ \norm{x}^2 = \norm{y}^2 $$
Wtedy:
\begin{equation}\label{eq:norm_of_real_vec}
    \norm{z} =\sqrt{\norm{x}^2 + \norm{y}^2 } = 1 \implies \norm{x}^2 = \norm{y}^2 = \frac{1}{2} 
\end{equation}
Wektory $x, y$, rzeczywiście, tworzą bazę ortogonalną.

Oznaczmy wartość własną jako $\lambda = a + ib$, gdzie $a, b \in \R$. Wyraźmy operator $A$ w bazie $\mathrm{lin}(x, y)$:
$$ Az = \lambda z, \quad Az = Ax + iAy  $$
$$ Ax + iAy = (a + ib)(x + iy) = ax + iay + ibx - by = (ax - by) + i(bx + ay) $$
Jeśli porównamy część rzeczywistą oraz część urojną, to mamy:
\begin{equation}\label{eq:operator_A_coeffs}
    Ax = ax - by, \quad Ay = bx + ay
\end{equation}
Dla wartości własnej sprzężonej wynik $\bar{\lambda}$ jest ten sam. W ten sposób przedstawiliśmy wektory $Ax, Ay$ za pomocą wektorów bazowych $x, y$. Jeśli potraktujemy współczynniki przy wektorach bazowych jako kolumny macierzy operatora $A$ to mamy:
\begin{equation}\label{eq:real_op_matrix}
\begin{pmatrix}
a & b \\
-b & a
\end{pmatrix}
\end{equation}

Wektory własne $z_1, \ldots, z_n$ z bazy $U$ możemy zapisać jak:
\begin{equation}
    z_{2k-1} = x_k + iy_k, \quad \bar{z}_{2k} = x_k - iy_k, \quad z_{i} = x_i,
\end{equation}
gdzie $(k = 1, 2, \ldots, q;\quad i = 2q + 1, \ldots, n)$ oraz $2q$ jest ilość wektorów zespolonych wraz ze sprzężeniem, oraz od $2q + 1$ do $n$ ilość wektorów rzeczywistych. Tym wektorom własnym odpowiadają wartości własne w ten sam sposób:
\begin{equation}\label{eq:eigenvalues}
    \lambda_{2k-1} = a_k + ib_k, \quad \bar{\lambda}_{2k} = a_k - ib_k, \quad \lambda_{i} = a_i
\end{equation}
gdzie $(k = 1, 2, \ldots, q;\quad i = 2q + 1, \ldots, n)$ oraz $a_k, b_k, a_i \in \R$.

Wyżej pokazaliśmy \eqref{eq:real_vec_orthogonal}, że dla dowolnego wektora $z_{k} = x_k \pm iy_k$, wektory rzeczywiste $x_k, y_k \in V$ są ortogonalne. Pokażemy teraz, że dla dowolnych dwóch wektorów $z_m = x_m + iy_m$, $z_j = x_j + iy_j$ z bazy ortonormalnej $U$, wektory rzeczywistę $x_m, y_m, x_j, y_j$ są ortogonalne między sobą, gdzie $m \neq j$ oraz $m, j \in \{1, 2, \ldots, n\}$.

Z założenia bazy ortonormalnej wektory $z_m, z_j, \bar{z}_m, \bar{z}_j$ są ortogonalne, więc możemy ułożyć następujący układ równań:
\[
\begin{cases}
\inner{z_m}{z_j} = \inner{x_m + iy_m}{x_j + iy_j} = 0 \\
\inner{z_m}{\bar{z}_j} = \inner{x_m + iy_m}{x_j - iy_j} = 0
\end{cases}
\]
Z tego wynika:
\[
\begin{cases}
\inner{z_m}{z_j} = \inner{x_m}{x_j} + \inner{y_m}{y_j} + i(\inner{y_m}{x_j} - \inner{x_m}{y_j}) = 0 \\
\inner{z_m}{\bar{z}_j} = \inner{x_m}{x_j} - \inner{y_m}{y_j} + i(\inner{y_m}{x_j} + \inner{x_m}{y_j}) = 0
\end{cases}
\]
Część rzeczywista i urojona muszą być zero. Porównajmy części rzeczywiste i urojne:

\[
\text{(Re)} \qquad
\begin{cases}
\inner{x_m}{x_j} + \inner{y_m}{y_j} = 0 \\
\inner{x_m}{x_j} - \inner{y_m}{y_j} = 0
\end{cases}
\implies \inner{x_m}{x_j} = 0, \quad \inner{y_m}{y_j} = 0
\]

\[
\text{(Im)} \qquad
\begin{cases}
\inner{y_m}{x_j} - \inner{x_m}{y_j} = 0 \\
\inner{y_m}{x_j} + \inner{x_m}{y_j} = 0
\end{cases}
\implies \inner{y_m}{x_j} = 0, \quad \inner{x_m}{y_j} = 0
\]
Więc dla wszystkich wektorów z bazy $U$:
$$z_{2k-1} = x_k + iy_k, \quad \bar{z}_{2k} = x_k - iy_k, \quad z_i = x_i \quad (k = 1, 2, \ldots, q; \quad i = 2q + 1, \ldots, n)$$
wektory rzeczywiste $x_k, y_k, x_i$ są ortogonalne między sobą. 

Z wektorów:
\begin{equation}\label{eq:real_base}
x_1, \ y_1, \ \ldots, \ x_k, \ y_k \ \ldots, \ x_q, \ y_q, \quad x_{2q+1}, \ \ldots, \ x_i, \ \ldots, x_n
\end{equation}
możemy utworzyć bazę ortogonalną w przestrzeni euklidesowej $V$. Jak wynika z \eqref{eq:norm_of_real_vec} wszystkie wektory możemy łatwo znormalizować, ponieważ każdy z nich ma tę samą normę. Po normalizacji mamy bazę ortonormalną.
 
Korzystając z \eqref{eq:operator_A_coeffs} oraz postaci wartości własnych \eqref{eq:eigenvalues}:
\begin{equation}\label{eq:op_A_in_real_base}
   Ax_k = a_k x_k - b_k y_k, \quad Ay_k = b_k x_k + a_k y_k, \quad Ax_i = a_i x_i, 
\end{equation}
gdzie $(k = 1, 2, \ldots, q; \ i = 2q + 1, \ldots, n)$.

W bazie \eqref{eq:real_base} operator $A$ ma macierz klatkowo-diagonalną:
\begin{equation}
\begin{Bmatrix}
\cblock{a_1}{b_1},
&
\ldots,
&
\cblock{a_q}{b_q},
&
a_{2q+1},
&
\ldots,
&
a_n
\end{Bmatrix}
\end{equation}

Więc mamy następujące twierdzenie:
\begin{theorem}\label{theorem:real_normal_form}
    W przestrzeni euklidesowej $V$ istnieje baza ortonormalna, w której operator normalny $A$ ma macierz klatkowo-diagonalną:
\begin{equation}
   A = O
\begin{Bmatrix}
\cblock{a_1}{b_1},
&
\ldots,
&
\cblock{a_q}{b_q},
&
a_{2q+1},
&
\ldots,
&
a_n
\end{Bmatrix}
O^{-1}, \quad (O^T = O^{-1}),
\end{equation}
gdzie $O$ jest macierzą ortogonalną złożoną z wektorów \eqref{eq:real_base}.
\end{theorem}

Teraz możemy łatwo wyprowadzić postać normalną dla macierzy antysymetrycznej oraz ortogonalnej.

\begin{theorem}
Operator antysymetryczny rzeczywisty $K$ w przestrzeni unitarnej $\V$ ma wartości własne w postaci czysto urojonej: $\lambda = ib$, gdzie $b \in \R$  
\end{theorem}
\begin{proof}
Z definicji operator antysymetryczny spełnia warunek $K = -K^*$, co dla rzeczywistego $K$ jest równoważne $K = -K^T$, gdyż $K^* = K^T$
Operator antysymetryczny jest operatorem normalnym rzeczywistym ($KK^* = K^*K = -K^2$).
Na mocy Twierdzenia Spektralnego ma bazę ortonormalną złożoną z wektorów własnych. 
Z $K = -K^T$ wynika oraz na mocy Lematu~\ref{lemma:same_eigenvector} (Wspólny wektor własny):
$$ Kv = \lambda v, \quad K^Tv = \bar{\lambda} v$$ 
$$ -Kv = \bar{\lambda} v \implies  Kv = -\bar{\lambda} v \implies \lambda v = -\bar{\lambda} v $$ 
Ponieważ $v \neq 0$:
$$ \lambda = -\bar{\lambda} \implies a + ib = -(a - ib) \implies a = 0,$$
gdzie $a, b \in \R$.

Więc wszystkie wartości własne mają postać $\lambda = ib$, gdzie $b \in \R$.
\end{proof}

Jak widzimy z Twierdzenia wyżej, operator antysymetryczny $K$ ma wartości czysto urojone, a więc w bazie \eqref{eq:real_base} dla \eqref{eq:op_A_in_real_base} musimy założyć:
$$ a_1 = a_2 = \ldots = a_q = a_{2q+1} = \ldots = a_n = 0 $$
Wtedy:
\begin{equation}
   Kx_k = - b_k y_k, \quad Ky_k = b_k x_k, \quad Kx_i = 0, 
\end{equation}
gdzie $(k = 1, 2, \ldots, q; \ i = 2q + 1, \ldots, n)$.


Na mocy Twierdzenia~\ref{theorem:real_normal_form} w przestrzeni euklidesowej \textbf{macierz antysymetryczna} operatora $K$ jest ortogonalnie podobna do macierzy klatkowo-diagonalnej:
\begin{equation}\label{eq:real_skew_matrix_normal_form}
    K =
        O
\begin{Bmatrix}
\cblock{0}{b_1},
&
\ldots,
&
\cblock{0}{b_q},
&
0,
&
\ldots,
&
0
\end{Bmatrix}
O^{-1}, \quad (O^T = O^{-1}),
\end{equation}
gdzie $O$ jest macierzą ortogonalną złożoną z wektorów \eqref{eq:real_base}.

Postać \eqref{eq:real_skew_matrix_normal_form} jest \textbf{formą normalną rzeczywistą dla macierzy antysymetrycznej}.

Operator ortogonalny $O$ spełnia warunki operatora unitarnego ($O^*O = I$, gdyż $O^* = O^T$), a więc na mocy Twierdzenia Spektralnego~\ref{theorem:unitary_spectral} ma wartości własne równe jedynce z modułu: $|\lambda| = 1, \ \lambda \in \C$, wtedy w bazie \eqref{eq:real_base} dla wzorów \eqref{eq:op_A_in_real_base} musimy założyć:
$$ |\lambda_k| = a_k^2 + b_k^2 = 1, \quad a_i = \pm1, $$
gdzie $\lambda_k$ wartość własna operatora $O$ oraz $(k = 1, 2, \ldots, q; \ i = 2q + 1, \ldots, n)$.

Każdą liczbę zespoloną $\lambda_k = a_k + ib_k, \ a_k, b_k \in \R$ ze wzoru Eulera można zapisać następująco:
$$ \lambda = \cos{\varphi_k} + i\sin{\varphi_k}, $$ gdzie $a_k = \cos{\varphi_k}, \ b_k = \sin{\varphi_k}, \ \varphi_k \in \R$.
Wtedy wzory \eqref{eq:op_A_in_real_base} możemy zapisać jako:
\begin{equation}
   Ox_k = \cos{\varphi_k} x_k - \sin{\varphi_k} y_k, \quad Oy_k = \sin{\varphi_k} x_k + \cos{\varphi_k} y_k, \quad Ox_i = \pm x_i, 
\end{equation}
gdzie $(k = 1, 2, \ldots, q; \ i = 2q + 1, \ldots, n)$.

Więc z powyższego oraz Twierdzenia~\ref{theorem:real_normal_form} w przestrzeni euklidesowej \textbf{macierz ortogonalna} operatora $O$ jest ortogonalnie podobna do postaci klatkowo-diagonalnej:
\begin{equation}\label{eq:orthogonal_matrix_normal_form}
    O =
        O_1
\begin{Bmatrix}
\cblock{\cos{\varphi_1}}{\sin{\varphi_1}},
&
\ldots,
&
\cblock{\cos{\varphi_q}}{\sin{\varphi_q}},
&
\pm 1,
&
\ldots,
&
\pm 1
\end{Bmatrix}
O_1^{-1}, \quad (O_1^T = O_1^{-1}),
\end{equation}
gdzie $O_1$ jest macierzą ortogonalną złożoną z wektorów \eqref{eq:real_base}.


Postać \eqref{eq:orthogonal_matrix_normal_form} jest \textbf{formą normalną rzeczywistą dla macierzy ortogonalnej}.
Klatki $\cblock{\cos{\varphi_q}}{\sin{\varphi_q}}$ macierzy $O$ reprezentują obroty wektrów przestrzni o kąt $\varphi$, a wartości $\pm1$ reprezentują odbicia wektorów przestrzni.